%% preambulo.tex — incluído por main.tex

%%%% Pacotes do template SBC
\usepackage{graphicx,url}%
\usepackage[T1]{fontenc}%
\usepackage[utf8]{inputenc}%
\usepackage[brazil]{babel}%   texto em português

%%%% Pacotes adicionais deste manuscrito
\usepackage{multirow}%
\usepackage{amsmath,amssymb,amsfonts}%
\usepackage{xcolor}%
\usepackage{textcomp}%
\usepackage{booktabs}%        \toprule, \midrule, \bottomrule
\usepackage{array}%           colunas de largura fixa (tipo L)
\usepackage{makecell}%        cabeçalhos de tabela em duas linhas
\usepackage{tabularx}%        tabelas que ocupam toda a largura do texto
\usepackage{etoolbox}%

%% Centraliza todas as tabelas e figuras (a classe Springer fazia isso sozinha)
\AtBeginEnvironment{table}{\centering}
\AtBeginEnvironment{figure}{\centering}

\sloppy

%% Coluna de texto que quebra linha e mantém alinhamento à esquerda
\newcolumntype{L}[1]{>{\raggedright\arraybackslash}p{#1}}
%% Coluna que estica até completar a largura do texto (alinhada à esquerda)
\newcolumntype{Y}{>{\raggedright\arraybackslash}X}

%% Fonte das tabelas em um único lugar: editar aqui muda todas de uma vez.
%% Aparece logo abaixo da tabela, em letra menor.
\newcommand{\fontetabela}{\par\smallskip{\small Fonte: Elaborada pelos
autores, a partir da Coorte de 100 Milhões de Brasileiros vinculada ao
SINASC.}}

%% Ajuste comum às tabelas de subgrupos (7 colunas)
\newcommand{\tabcompacta}{\small\setlength{\tabcolsep}{4pt}}

\usepackage{tikz}%           fluxograma da metodologia
\usetikzlibrary{positioning,arrows.meta}
\usepackage{placeins}%       \FloatBarrier: impede tabelas de passarem para a seção seguinte
\usepackage{longtable}%       tabela que atravessa páginas (trabalhos relacionados)

\graphicspath{{figuras/}}% todas as imagens ficam em figuras/

\raggedbottom

%% Permite quebrar endereços longos (DOIs nas referências) em qualquer caractere
\makeatletter
\g@addto@macro{\UrlBreaks}{\do\a\do\b\do\c\do\d\do\e\do\f\do\g\do\h\do\i\do\j\do\k\do\l\do\m\do\n\do\o\do\p\do\q\do\r\do\s\do\t\do\u\do\v\do\w\do\x\do\y\do\z\do\A\do\B\do\C\do\D\do\E\do\F\do\G\do\H\do\I\do\J\do\K\do\L\do\M\do\N\do\O\do\P\do\Q\do\R\do\S\do\T\do\U\do\V\do\W\do\X\do\Y\do\Z\do0\do1\do2\do3\do4\do5\do6\do7\do8\do9}
\makeatother
